% !PRINTED 769
\weberpage{769}
%<<CONT

and consequently, if $u$, $v$ are indeterminate quantities and the
\gesperrt{congruence is referred only to the coefficients},
\[
(u+v)^p \equiv u^p+v^p \ (\mathrm{mod}\ p).
\]

If we again replace $v$ in it by a sum $v+w+\cdots$,
we obtain by complete induction (or else from the
polynomial theorem):
\[
(u+v+w+\cdots)^p \equiv u^p+v^p+w^p+\cdots \ (\mathrm{mod}\ p),
\]
and if we apply this formula repeatedly, and understand by $k$
an arbitrary positive exponent,
\weq{1}{(u+v+w+\cdots)^{p^k} \equiv u^{p^k}+v^{p^k}+w^{p^k}+\cdots \ (\mathrm{mod}\ p).}

Let $x$ now be a variable, and let $a_1, a_2, \dots, a_\nu$ be integral rational
numbers,
\[
f(x) = x^\nu + a_1 x^{\nu-1} + a_2 x^{\nu-2} + \cdots + a_\nu
\]
an integral function of $x$, whose roots we may take to be $\alpha, \beta, \gamma, \dots$,
so that
\[
-a_1 = \Sigma\alpha, \quad a_2 = \Sigma\alpha\beta, \quad -a_3 = \Sigma\alpha\beta\gamma \dots
\]
are the symmetric fundamental functions of the $\alpha, \beta, \gamma, \dots$.

We form the function
\[
F(x) = x^\nu + A_1 x^{\nu-1} + A_2 x^{\nu-2} + \cdots + A_\nu,
\]
whose roots are the $p^{k}$th powers of the roots of $f$, that is,
\[
\alpha^{p^k}, \ \beta^{p^k}, \ \gamma^{p^k}, \dots
\]
The coefficients of $F$:
\[
A_1 = -\Sigma\alpha^{p^k}, \quad A_2 = \Sigma\alpha^{p^k}\beta^{p^k}, \quad A_3 = -\Sigma\alpha^{p^k}\beta^{p^k}\gamma^{p^k} \dots
\]
are, as symmetric functions of the $\alpha, \beta, \gamma, \dots$ with integral
coefficients, representable rationally and with integral coefficients by means of the $a_1, a_2, \dots$, and are
therefore integers. The quotients
\[
\frac{A_1 - a_1^{p^k}}{p}, \quad \frac{A_2 - a_2^{p^k}}{p}, \quad \frac{A_3 - a_3^{p^k}}{p}, \dots
\]
are, however, by formula (1) likewise integers, and thus by
\gesperrt{Fermat's theorem}
\weq{2}{F(x) \equiv f(x) \ (\mathrm{mod}\ p).}

These theorems are now to be applied to the function $f_n(x)$,
whose roots are all the \gesperrt{primitive} $n$th roots of unity $\alpha, \beta, \gamma, \dots$, and which, as was proved in \S.\,133 of the first
volume, has integral coefficients $a_1, a_2, \dots$.
